\documentclass[11pt,a4paper]{article}
\usepackage[utf8]{inputenc}
\usepackage[T1]{fontenc}
\usepackage[spanish,es-minimal,es-noshorthands]{babel}
\usepackage{lmodern}
\usepackage[sfdefault,lining]{FiraSans}
\usepackage[a4paper,top=2.3cm,bottom=2.2cm,left=1.8cm,right=1.8cm,headheight=1.3cm,headsep=0.4cm,footskip=1cm]{geometry}
\usepackage{amsmath,amssymb}
\usepackage{icomma}
\usepackage{graphicx}
\usepackage[table]{xcolor}
\usepackage{tikz}
\usepackage{fontawesome5}
\usepackage{ccicons}
\usepackage{fancyhdr}
\usepackage{lastpage}
\usepackage{tasks}
\usepackage{needspace}
\usepackage[most]{tcolorbox}
\usepackage{array}
\usepackage[hidelinks]{hyperref}
\AtBeginDocument{\hypersetup{pdftitle={Ficha de ejercicios de la unidad 2: Números reales. 3.º de ESO},pdfsubject={Matemáticas 3.º ESO},pdfkeywords={números reales, potencias, radicales, notación científica, intervalos}}}

\definecolor{navy}{HTML}{1B2748}
\definecolor{accent}{HTML}{E8603A}
\definecolor{teal}{HTML}{1F9E90}
\definecolor{soft}{HTML}{EEF1F8}
\definecolor{muted}{HTML}{5D6685}

\setlength{\parindent}{0pt}
\setlength{\parskip}{3pt}
\settasks{label=\alph*), label-format=\bfseries\color{navy}, label-width=1.3em, item-indent=1.8em, column-sep=0.6em, after-item-skip=2pt, before-skip=3pt, after-skip=4pt}

% ---------- capçalera i peu ----------
\newcommand{\llicencia}{\href{https://creativecommons.org/licenses/by-nc-sa/4.0/deed.es}{\ccbyncsa}\ \ Licencia Creative Commons Atribución-NoComercial-CompartirIgual 4.0 Internacional (CC BY-NC-SA 4.0)}
\pagestyle{fancy}
\fancyhf{}
\fancyhead[L]{\raisebox{-0.35\height}{\includegraphics[height=0.95cm]{logo-gva.png}}}
\fancyhead[C]{\small\color{muted}Ficha de ejercicios · Unidad 2: Números reales · 3.º de ESO}
\fancyhead[R]{\raisebox{-0.35\height}{\includegraphics[height=0.95cm]{logo-vermellar.png}}}
\fancyfoot[L]{\scriptsize\color{muted}\llicencia}
\fancyfoot[R]{\small\color{muted}\thepage}
\renewcommand{\headrulewidth}{0.4pt}
\renewcommand{\headrule}{\color{muted!40}\hrule height \headrulewidth}
\fancypagestyle{primera}{\fancyhf{}\renewcommand{\headrulewidth}{0pt}
  \fancyfoot[L]{\scriptsize\color{muted}\llicencia}
  \fancyfoot[R]{\small\color{muted}\thepage}}

% ---------- nivells i exercicis ----------
\newcommand{\estrelles}[1]{{\color{accent}\ifcase#1\or\faStar\faStar[regular]\faStar[regular]\or\faStar\faStar\faStar[regular]\or\faStar\faStar\faStar\fi}}
\newcounter{exer}
\newcommand{\ex}[1]{\par\needspace{4\baselineskip}\medskip\stepcounter{exer}\noindent{\color{navy}\bfseries\theexer.}\ \ {\small\estrelles{#1}}\ \ }

\newcommand{\seccio}[2]{\par\needspace{10\baselineskip}\bigskip
  \begin{tcolorbox}[colback=navy,colframe=navy,coltext=white,arc=2mm,boxrule=0pt,left=3mm,right=3mm,top=1.2mm,bottom=1.2mm]
  {\large\bfseries #1\quad #2}
  \end{tcolorbox}\vspace{-2pt}}

% ---------- recta real ----------
\newcommand{\rectabuida}{\begin{tikzpicture}[x=0.42cm,baseline=0]
  \draw[-latex,muted] (-6.7,0)--(6.8,0);
  \foreach \i in {-6,...,6}{\draw[muted] (\i,0.12)--(\i,-0.12) node[below=-1pt,font=\scriptsize,text=muted]{$\i$};}
  \end{tikzpicture}}
% \rectaint{x1}{x2}{estil esquerra}{estil dreta}: estil = ple, buit, inf
\newcommand{\extrem}[2]{\ifx#2p\fill[accent] (#1,0) circle (0.13);\else\ifx#2b\filldraw[fill=white,draw=accent,line width=1.2pt] (#1,0) circle (0.13);\fi\fi}
\newcommand{\rectaint}[4]{\begin{tikzpicture}[x=0.42cm,baseline=0]
  \draw[-latex,muted] (-6.7,0)--(6.8,0);
  \foreach \i in {-6,...,6}{\draw[muted] (\i,0.12)--(\i,-0.12) node[below=-1pt,font=\scriptsize,text=muted]{$\i$};}
  \draw[accent,line width=3pt] (#1,0)--(#2,0);
  \ifx#3i\fill[accent] (#1-0.25,0)--(#1+0.15,0.2)--(#1+0.15,-0.2)--cycle;\else\extrem{#1}{#3}\fi
  \ifx#4i\fill[accent] (#2+0.25,0)--(#2-0.15,0.2)--(#2-0.15,-0.2)--cycle;\else\extrem{#2}{#4}\fi
  \end{tikzpicture}}

\newcolumntype{C}[1]{>{\centering\arraybackslash}p{#1}}

% claudàtors oberts dels intervals: ]a,b[ (amb l'espaiat de parèntesi, perquè ]-3 no siga una resta)
\newcommand{\lo}{\mathopen{]}}
\newcommand{\ro}{\mathclose{[}}
% =====================================================================
% Capçalera FentMates de les fitxes: un color per a cada curs.
% Els colors són els de la pàgina inicial (LEVELS a assets/js/materials.js);
% el logo ve de _eines/logo/pdf/ (python3 _eines/logo/genera.py).
%
% Ús, al final del preàmbul de la unitat:
%   % =====================================================================
% Capçalera FentMates de les fitxes: un color per a cada curs.
% Els colors són els de la pàgina inicial (LEVELS a assets/js/materials.js);
% el logo ve de _eines/logo/pdf/ (python3 _eines/logo/genera.py).
%
% Ús, al final del preàmbul de la unitat:
%   % =====================================================================
% Capçalera FentMates de les fitxes: un color per a cada curs.
% Els colors són els de la pàgina inicial (LEVELS a assets/js/materials.js);
% el logo ve de _eines/logo/pdf/ (python3 _eines/logo/genera.py).
%
% Ús, al final del preàmbul de la unitat:
%   \input{../comu/fentmates.tex}
%   \fmcurs{1eso}
% i, al començament del document:
%   \capcalera{Matemàtiques}{1r ESO}{Fitxa d'exercicis · Unitat 2}{Divisibilitat}
% Opcional abans de \input: \def\fmneed{14} (espai mínim abans de cada secció).
% =====================================================================
\makeatletter
\def\fm@logodir{../../logo/pdf}
\providecommand{\fmneed}{10}
\def\fm@cursnom{}

% clau del curs, color i color del text damunt del color
\newcommand{\fm@curs}[3]{\@namedef{fm@color@#1}{#2}\@namedef{fm@text@#1}{#3}}
\fm@curs{1eso}{E6B84A}{navy}
\fm@curs{2eso}{3B82F6}{white}
\fm@curs{3eso}{FF7A50}{white}
\fm@curs{4eso-a}{A78BFA}{white}
\fm@curs{4eso-b}{6D28D9}{white}
\fm@curs{1batx-mi}{9B1C1C}{white}
\fm@curs{1batx-ccss}{EF6B6B}{white}
\fm@curs{2batx-mii}{1E6B47}{white}
\fm@curs{2batx-ccss}{2BB3A3}{white}

\newcommand{\fmcurs}[1]{%
  \@ifundefined{fm@color@#1}{\PackageError{fentmates}{Curs desconegut: #1}{}}{}%
  \edef\fm@tmp{\noexpand\definecolor{curs}{HTML}{\@nameuse{fm@color@#1}}%
               \noexpand\colorlet{curstext}{\@nameuse{fm@text@#1}}}\fm@tmp
  \colorlet{cursfosc}{curs!72!black}%
  \def\fm@logo{\fm@logodir/vertical-#1.pdf}%
  \def\fm@marca{\fm@logodir/marca-#1.pdf}}

% ---------- capçalera de la primera pàgina ----------
% #1 matèria (petit), #2 curs, #3 tipus de document i unitat, #4 títol
\newcommand{\capcalera}[4]{%
  \gdef\fm@cursnom{#2}%
  \par\noindent
  \begin{tcolorbox}[colback=curs,colframe=curs,boxrule=0pt,arc=2.2mm,width=3.5cm,height=2.15cm,
      valign=center,left=3mm,right=2mm,top=1mm,bottom=1mm,nobeforeafter,box align=center]
    \color{curstext}{\footnotesize #1}\par\vspace{1pt}{\fontsize{19}{21}\selectfont\bfseries #2\par}%
  \end{tcolorbox}\hspace{5mm}%
  \begin{minipage}[c]{\dimexpr\linewidth-3.5cm-5mm-2.9cm\relax}
    {\color{cursfosc}\small\bfseries\MakeUppercase{#3}}\par\vspace{2pt}
    {\color{navy}\fontsize{20}{23}\selectfont\bfseries #4\par}
  \end{minipage}\hfill
  \begin{minipage}[c]{2.6cm}\raggedleft\includegraphics[width=2.6cm]{\fm@logo}\end{minipage}\par
  \vspace{2pt}}

% ---------- barres de secció del color del curs ----------
\renewcommand{\seccio}[2]{\par\needspace{\fmneed\baselineskip}\bigskip
  \begin{tcolorbox}[colback=curs,colframe=curs,coltext=curstext,arc=2mm,boxrule=0pt,left=3mm,right=3mm,top=1.2mm,bottom=1.2mm]
  {\large\bfseries #1\quad #2}
  \end{tcolorbox}\vspace{-2pt}}

% ---------- estrelles de dificultat del color del curs (a Batxillerat no se'n posen) ----------
\@ifundefined{estrelles}{}{%
  \renewcommand{\estrelles}[1]{{\color{curs}\ifcase#1\or\faStar\faStar[regular]\faStar[regular]\or\faStar\faStar\faStar[regular]\or\faStar\faStar\faStar\fi}}}

% ---------- capçalera i peu de les pàgines ----------
% primera pàgina: logos institucionals al peu
\fancypagestyle{primera}{\fancyhf{}\renewcommand{\headrulewidth}{0pt}
  \fancyfoot[L]{\raisebox{-0.32\height}{\includegraphics[height=0.78cm]{logo-gva.png}}\hspace{5mm}%
    \raisebox{-0.32\height}{\includegraphics[height=0.78cm]{logo-vermellar.png}}}
  \fancyfoot[R]{\scriptsize\color{muted}\href{https://creativecommons.org/licenses/by-nc-sa/4.0/deed.ca}{\ccbyncsa}\ \ CC BY-NC-SA 4.0 · fentmates.org\qquad\small\thepage}}
% la resta: el símbol a l'esquerra i el curs a la dreta, amb el títol al centre com abans
\fancyhead[L]{\raisebox{-0.3\height}{\includegraphics[height=0.62cm]{\fm@marca}}}
\fancyhead[R]{\small\bfseries\color{cursfosc}\fm@cursnom}
\renewcommand{\headrule}{\color{curs!45}\hrule height \headrulewidth}
\makeatother

%   \fmcurs{1eso}
% i, al començament del document:
%   \capcalera{Matemàtiques}{1r ESO}{Fitxa d'exercicis · Unitat 2}{Divisibilitat}
% Opcional abans de \input: \def\fmneed{14} (espai mínim abans de cada secció).
% =====================================================================
\makeatletter
\def\fm@logodir{../../logo/pdf}
\providecommand{\fmneed}{10}
\def\fm@cursnom{}

% clau del curs, color i color del text damunt del color
\newcommand{\fm@curs}[3]{\@namedef{fm@color@#1}{#2}\@namedef{fm@text@#1}{#3}}
\fm@curs{1eso}{E6B84A}{navy}
\fm@curs{2eso}{3B82F6}{white}
\fm@curs{3eso}{FF7A50}{white}
\fm@curs{4eso-a}{A78BFA}{white}
\fm@curs{4eso-b}{6D28D9}{white}
\fm@curs{1batx-mi}{9B1C1C}{white}
\fm@curs{1batx-ccss}{EF6B6B}{white}
\fm@curs{2batx-mii}{1E6B47}{white}
\fm@curs{2batx-ccss}{2BB3A3}{white}

\newcommand{\fmcurs}[1]{%
  \@ifundefined{fm@color@#1}{\PackageError{fentmates}{Curs desconegut: #1}{}}{}%
  \edef\fm@tmp{\noexpand\definecolor{curs}{HTML}{\@nameuse{fm@color@#1}}%
               \noexpand\colorlet{curstext}{\@nameuse{fm@text@#1}}}\fm@tmp
  \colorlet{cursfosc}{curs!72!black}%
  \def\fm@logo{\fm@logodir/vertical-#1.pdf}%
  \def\fm@marca{\fm@logodir/marca-#1.pdf}}

% ---------- capçalera de la primera pàgina ----------
% #1 matèria (petit), #2 curs, #3 tipus de document i unitat, #4 títol
\newcommand{\capcalera}[4]{%
  \gdef\fm@cursnom{#2}%
  \par\noindent
  \begin{tcolorbox}[colback=curs,colframe=curs,boxrule=0pt,arc=2.2mm,width=3.5cm,height=2.15cm,
      valign=center,left=3mm,right=2mm,top=1mm,bottom=1mm,nobeforeafter,box align=center]
    \color{curstext}{\footnotesize #1}\par\vspace{1pt}{\fontsize{19}{21}\selectfont\bfseries #2\par}%
  \end{tcolorbox}\hspace{5mm}%
  \begin{minipage}[c]{\dimexpr\linewidth-3.5cm-5mm-2.9cm\relax}
    {\color{cursfosc}\small\bfseries\MakeUppercase{#3}}\par\vspace{2pt}
    {\color{navy}\fontsize{20}{23}\selectfont\bfseries #4\par}
  \end{minipage}\hfill
  \begin{minipage}[c]{2.6cm}\raggedleft\includegraphics[width=2.6cm]{\fm@logo}\end{minipage}\par
  \vspace{2pt}}

% ---------- barres de secció del color del curs ----------
\renewcommand{\seccio}[2]{\par\needspace{\fmneed\baselineskip}\bigskip
  \begin{tcolorbox}[colback=curs,colframe=curs,coltext=curstext,arc=2mm,boxrule=0pt,left=3mm,right=3mm,top=1.2mm,bottom=1.2mm]
  {\large\bfseries #1\quad #2}
  \end{tcolorbox}\vspace{-2pt}}

% ---------- estrelles de dificultat del color del curs (a Batxillerat no se'n posen) ----------
\@ifundefined{estrelles}{}{%
  \renewcommand{\estrelles}[1]{{\color{curs}\ifcase#1\or\faStar\faStar[regular]\faStar[regular]\or\faStar\faStar\faStar[regular]\or\faStar\faStar\faStar\fi}}}

% ---------- capçalera i peu de les pàgines ----------
% primera pàgina: logos institucionals al peu
\fancypagestyle{primera}{\fancyhf{}\renewcommand{\headrulewidth}{0pt}
  \fancyfoot[L]{\raisebox{-0.32\height}{\includegraphics[height=0.78cm]{logo-gva.png}}\hspace{5mm}%
    \raisebox{-0.32\height}{\includegraphics[height=0.78cm]{logo-vermellar.png}}}
  \fancyfoot[R]{\scriptsize\color{muted}\href{https://creativecommons.org/licenses/by-nc-sa/4.0/deed.ca}{\ccbyncsa}\ \ CC BY-NC-SA 4.0 · fentmates.org\qquad\small\thepage}}
% la resta: el símbol a l'esquerra i el curs a la dreta, amb el títol al centre com abans
\fancyhead[L]{\raisebox{-0.3\height}{\includegraphics[height=0.62cm]{\fm@marca}}}
\fancyhead[R]{\small\bfseries\color{cursfosc}\fm@cursnom}
\renewcommand{\headrule}{\color{curs!45}\hrule height \headrulewidth}
\makeatother

%   \fmcurs{1eso}
% i, al començament del document:
%   \capcalera{Matemàtiques}{1r ESO}{Fitxa d'exercicis · Unitat 2}{Divisibilitat}
% Opcional abans de \input: \def\fmneed{14} (espai mínim abans de cada secció).
% =====================================================================
\makeatletter
\def\fm@logodir{../../logo/pdf}
\providecommand{\fmneed}{10}
\def\fm@cursnom{}

% clau del curs, color i color del text damunt del color
\newcommand{\fm@curs}[3]{\@namedef{fm@color@#1}{#2}\@namedef{fm@text@#1}{#3}}
\fm@curs{1eso}{E6B84A}{navy}
\fm@curs{2eso}{3B82F6}{white}
\fm@curs{3eso}{FF7A50}{white}
\fm@curs{4eso-a}{A78BFA}{white}
\fm@curs{4eso-b}{6D28D9}{white}
\fm@curs{1batx-mi}{9B1C1C}{white}
\fm@curs{1batx-ccss}{EF6B6B}{white}
\fm@curs{2batx-mii}{1E6B47}{white}
\fm@curs{2batx-ccss}{2BB3A3}{white}

\newcommand{\fmcurs}[1]{%
  \@ifundefined{fm@color@#1}{\PackageError{fentmates}{Curs desconegut: #1}{}}{}%
  \edef\fm@tmp{\noexpand\definecolor{curs}{HTML}{\@nameuse{fm@color@#1}}%
               \noexpand\colorlet{curstext}{\@nameuse{fm@text@#1}}}\fm@tmp
  \colorlet{cursfosc}{curs!72!black}%
  \def\fm@logo{\fm@logodir/vertical-#1.pdf}%
  \def\fm@marca{\fm@logodir/marca-#1.pdf}}

% ---------- capçalera de la primera pàgina ----------
% #1 matèria (petit), #2 curs, #3 tipus de document i unitat, #4 títol
\newcommand{\capcalera}[4]{%
  \gdef\fm@cursnom{#2}%
  \par\noindent
  \begin{tcolorbox}[colback=curs,colframe=curs,boxrule=0pt,arc=2.2mm,width=3.5cm,height=2.15cm,
      valign=center,left=3mm,right=2mm,top=1mm,bottom=1mm,nobeforeafter,box align=center]
    \color{curstext}{\footnotesize #1}\par\vspace{1pt}{\fontsize{19}{21}\selectfont\bfseries #2\par}%
  \end{tcolorbox}\hspace{5mm}%
  \begin{minipage}[c]{\dimexpr\linewidth-3.5cm-5mm-2.9cm\relax}
    {\color{cursfosc}\small\bfseries\MakeUppercase{#3}}\par\vspace{2pt}
    {\color{navy}\fontsize{20}{23}\selectfont\bfseries #4\par}
  \end{minipage}\hfill
  \begin{minipage}[c]{2.6cm}\raggedleft\includegraphics[width=2.6cm]{\fm@logo}\end{minipage}\par
  \vspace{2pt}}

% ---------- barres de secció del color del curs ----------
\renewcommand{\seccio}[2]{\par\needspace{\fmneed\baselineskip}\bigskip
  \begin{tcolorbox}[colback=curs,colframe=curs,coltext=curstext,arc=2mm,boxrule=0pt,left=3mm,right=3mm,top=1.2mm,bottom=1.2mm]
  {\large\bfseries #1\quad #2}
  \end{tcolorbox}\vspace{-2pt}}

% ---------- estrelles de dificultat del color del curs (a Batxillerat no se'n posen) ----------
\@ifundefined{estrelles}{}{%
  \renewcommand{\estrelles}[1]{{\color{curs}\ifcase#1\or\faStar\faStar[regular]\faStar[regular]\or\faStar\faStar\faStar[regular]\or\faStar\faStar\faStar\fi}}}

% ---------- capçalera i peu de les pàgines ----------
% primera pàgina: logos institucionals al peu
\fancypagestyle{primera}{\fancyhf{}\renewcommand{\headrulewidth}{0pt}
  \fancyfoot[L]{\raisebox{-0.32\height}{\includegraphics[height=0.78cm]{logo-gva.png}}\hspace{5mm}%
    \raisebox{-0.32\height}{\includegraphics[height=0.78cm]{logo-vermellar.png}}}
  \fancyfoot[R]{\scriptsize\color{muted}\href{https://creativecommons.org/licenses/by-nc-sa/4.0/deed.ca}{\ccbyncsa}\ \ CC BY-NC-SA 4.0 · fentmates.org\qquad\small\thepage}}
% la resta: el símbol a l'esquerra i el curs a la dreta, amb el títol al centre com abans
\fancyhead[L]{\raisebox{-0.3\height}{\includegraphics[height=0.62cm]{\fm@marca}}}
\fancyhead[R]{\small\bfseries\color{cursfosc}\fm@cursnom}
\renewcommand{\headrule}{\color{curs!45}\hrule height \headrulewidth}
\makeatother

\fmcurs{3eso}

\begin{document}
\thispagestyle{primera}

% ---------- encapçalament ----------
\capcalera{Matemáticas}{3.º ESO}{Ficha de ejercicios · Unidad 2}{Números reales}
\vspace{10pt}
Nombre y apellidos: \rule{6.6cm}{0.4pt}\hfill Grupo: \rule{1.6cm}{0.4pt}\hfill Fecha: \rule{2.2cm}{0.4pt}
\vspace{8pt}

\begin{tcolorbox}[colback=soft,colframe=soft,arc=2mm,left=3mm,right=3mm,top=1.5mm,bottom=1.5mm]
\small Ejercicios graduados por niveles de dificultad:\quad
\estrelles{1}\ básico \qquad \estrelles{2}\ medio \qquad \estrelles{3}\ avanzado.\\[2pt]
Resuélvelos en el cuaderno mostrando todos los pasos. Las rectas y las tablas se pueden completar directamente en la ficha.
\end{tcolorbox}

% =====================================================
\seccio{1}{Potencias de exponente natural}

\ex{1} Escribe como potencia y calcula:
\begin{tasks}(2)
  \task $5\cdot5\cdot5$
  \task $(-2)(-2)(-2)(-2)(-2)$
  \task $\frac23\cdot\frac23\cdot\frac23\cdot\frac23$
  \task $(-1)(-1)(-1)(-1)(-1)(-1)(-1)$
\end{tasks}

\ex{1} Calcula:
\begin{tasks}(4)
  \task $2^{6}$
  \task $(-3)^{3}$
  \task $(-3)^{4}$
  \task $-3^{4}$
  \task $(-1)^{100}$
  \task $\left(-\frac12\right)^{3}$
  \task $0^{9}$
  \task $10^{7}$
\end{tasks}

\ex{2} Sin hacer ningún cálculo, indica si el resultado es positivo o negativo. Justifícalo.
\begin{tasks}(4)
  \task $(-7)^{12}$
  \task $(-5)^{9}$
  \task $-(-2)^{6}$
  \task $-(-4)^{3}$
\end{tasks}

\ex{2} Calcula:
\begin{tasks}(2)
  \task $-2^{2}+(-2)^{2}-(-2)^{3}$
  \task $(-1)^{5}+(-1)^{6}-1^{7}$
  \task $3^{2}-2^{3}+(-1)^{10}$
  \task $\left(-\frac23\right)^{2}\cdot\left(\frac32\right)^{3}$
\end{tasks}

\ex{3} Halla el valor de $x$ en cada caso:
\begin{tasks}(4)
  \task $2^{x}=64$
  \task $(-3)^{x}=-27$
  \task $x^{4}=81$
  \task $\left(\frac25\right)^{x}=\frac{8}{125}$
\end{tasks}

% =====================================================
\seccio{2}{Propiedades de las potencias}

\ex{1} Expresa como una sola potencia:
\begin{tasks}(3)
  \task $3^{4}\cdot3^{5}$
  \task $7^{10}:7^{6}$
  \task $\left(2^{3}\right)^{5}$
  \task $5^{3}\cdot5\cdot5^{4}$
  \task $a^{8}:a^{3}\ \ (a\neq0)$
  \task $\left((-2)^{2}\right)^{3}$
\end{tasks}

\ex{1} Aplica las propiedades de las potencias con el mismo exponente y calcula:
\begin{tasks}(4)
  \task $2^{5}\cdot5^{5}$
  \task $12^{4}:4^{4}$
  \task $(-3)^{3}\cdot(-2)^{3}$
  \task $\left(\frac34\right)^{2}\cdot\left(\frac43\right)^{2}$
\end{tasks}

\ex{2} Simplifica y calcula:
\begin{tasks}(4)
  \task $\dfrac{2^{7}\cdot2^{3}}{\left(2^{4}\right)^{2}}$
  \task $\dfrac{\left(3^{2}\right)^{4}\cdot3^{3}}{3^{5}\cdot3^{4}}$
  \task $\dfrac{5^{6}\cdot5^{3}}{\left(5^{2}\right)^{4}}$
  \task $\dfrac{6^{3}\cdot6^{5}}{\left(6^{2}\right)^{3}}$
\end{tasks}

\ex{2} Descompón las bases en factores primos, simplifica y calcula:
\begin{tasks}(4)
  \task $\dfrac{4^{3}\cdot8^{2}}{16^{2}}$
  \task $\dfrac{9^{2}\cdot27}{3^{6}}$
  \task $\dfrac{25^{3}}{5^{2}\cdot125}$
  \task $\dfrac{6^{4}}{2^{3}\cdot3^{2}}$
\end{tasks}

\ex{3} Descompón, simplifica y calcula:
\begin{tasks}(4)
  \task $\dfrac{12^{3}\cdot18^{2}}{8^{2}\cdot27^{2}}$
  \task $\dfrac{15^{2}\cdot10^{3}}{6^{2}\cdot25^{2}}$
  \task $\dfrac{20^{3}\cdot9^{2}}{30^{2}\cdot4^{2}\cdot3}$
  \task $\dfrac{14^{2}\cdot21^{3}}{49^{2}\cdot6^{2}}$
\end{tasks}

\ex{3} Indica si cada igualdad es verdadera o falsa. Justifica la respuesta.
\begin{tasks}(3)
  \task $2^{3}+2^{3}=2^{4}$
  \task $3^{2}\cdot3^{3}=9^{5}$
  \task $(2+3)^{2}=2^{2}+3^{2}$
  \task $5^{3}\cdot2^{3}=10^{3}$
  \task $2^{10}+2^{10}=2^{11}$
  \task $4^{5}:4^{5}=0$
\end{tasks}

% =====================================================
\seccio{3}{Potencias de exponente entero negativo}

\ex{1} Calcula:
\begin{tasks}(4)
  \task $3^{-2}$
  \task $(-2)^{-3}$
  \task $10^{-3}$
  \task $\left(\frac14\right)^{-2}$
  \task $\left(\frac23\right)^{-3}$
  \task $\left(-\frac35\right)^{-2}$
  \task $7^{0}$
  \task $-2^{-4}$
\end{tasks}

\ex{1} Escribe como una potencia de exponente negativo:
\begin{tasks}(5)
  \task $\dfrac{1}{5^{4}}$
  \task $\dfrac{1}{27}$
  \task $\dfrac{1}{16}$
  \task $0,001$
  \task $\dfrac{1}{49}$
\end{tasks}

\ex{2} Expresa como una sola potencia:
\begin{tasks}(3)
  \task $5^{-3}\cdot5^{7}$
  \task $2^{-4}:2^{3}$
  \task $\left(3^{-2}\right)^{-3}$
  \task $7^{-2}\cdot7^{-5}:7^{-9}$
  \task $\left(4^{-1}\right)^{3}\cdot4^{5}$
  \task $10^{-3}\cdot10^{5}\cdot10^{-6}$
\end{tasks}

\ex{2} Calcula y da el resultado en forma de fracción:
\begin{tasks}(4)
  \task $2^{-1}+2^{-2}$
  \task $3^{-1}-3^{-2}$
  \task $\left(\frac12\right)^{-2}+\left(\frac13\right)^{-1}$
  \task $5^{0}-5^{-1}$
\end{tasks}

\ex{3} Descompón, simplifica y calcula:
\begin{tasks}(3)
  \task $\dfrac{8^{-2}\cdot4^{3}}{2^{-5}\cdot16}$
  \task $\dfrac{9^{-1}\cdot27^{2}}{3^{-2}\cdot81}$
  \task $\dfrac{6^{-2}\cdot12^{3}}{18^{-1}\cdot2^{5}}$
  \task $\dfrac{10^{-2}\cdot25^{3}}{5^{4}\cdot2^{-3}}$
  \task $\left(\dfrac23\right)^{-2}\cdot\left(\dfrac94\right)^{-1}$
  \task $\left[\left(\dfrac12\right)^{-3}\cdot4^{-2}\right]^{-1}$
\end{tasks}

% =====================================================
\seccio{4}{Notación científica}

\ex{1} Escribe en notación científica:
\begin{tasks}(3)
  \task $45\,000\,000$
  \task $0,000\,72$
  \task $908\,000$
  \task $0,000\,000\,3$
  \task $-12\,500$
  \task $0,05$
\end{tasks}

\ex{1} Escribe en forma decimal:
\begin{tasks}(5)
  \task $3,2\cdot10^{6}$
  \task $5\cdot10^{-3}$
  \task $1,07\cdot10^{9}$
  \task $8,45\cdot10^{-6}$
  \task $-6,1\cdot10^{4}$
\end{tasks}

\ex{2} Indica cuáles de estas expresiones no están escritas en notación científica y corrígelas:
\begin{tasks}(3)
  \task $2,7\cdot10^{5}$
  \task $34\cdot10^{6}$
  \task $0,8\cdot10^{-3}$
  \task $10,5\cdot10^{-2}$
  \task $125\cdot10^{-8}$
  \task $0,0042\cdot10^{9}$
\end{tasks}

\ex{2} Ordena estos números de menor a mayor:
\[
3,1\cdot10^{-4}\qquad 2,9\cdot10^{-3}\qquad 5\cdot10^{-5}\qquad 1,2\cdot10^{-4}\qquad 9\cdot10^{-6}
\]

\ex{3} Expresa en notación científica los siguientes datos:
\begin{tasks}(1)
  \task La velocidad de la luz es de $300\,000$ km/s.
  \task El diámetro de un glóbulo rojo es de $0,000\,007\,5$ m.
  \task La población mundial es de unos 8\,000 millones de personas.
  \task La masa de un átomo de hidrógeno es de $0,000\,000\,000\,000\,000\,000\,000\,001\,67$ g.
  \task Un año luz equivale a unos $9\,460\,000\,000\,000$ km.
\end{tasks}

% =====================================================
\seccio{5}{Radicales. Extracción e introducción de factores}

\ex{1} Calcula, si existen en $\mathbb{R}$:
\begin{tasks}(4)
  \task $\sqrt{81}$
  \task $\sqrt[3]{-125}$
  \task $\sqrt[4]{16}$
  \task $\sqrt{-9}$
  \task $\sqrt[5]{-32}$
  \task $\sqrt{\dfrac{25}{49}}$
  \task $\sqrt[3]{\dfrac{8}{27}}$
  \task $\sqrt[6]{1}$
\end{tasks}

\ex{1} Extrae todos los factores posibles del radical:
\begin{tasks}(6)
  \task $\sqrt{12}$
  \task $\sqrt{18}$
  \task $\sqrt{45}$
  \task $\sqrt{98}$
  \task $\sqrt[3]{16}$
  \task $\sqrt[3]{81}$
\end{tasks}

\ex{2} Extrae todos los factores posibles del radical:
\begin{tasks}(3)
  \task $\sqrt{200}$
  \task $\sqrt{2^{5}\cdot3^{3}}$
  \task $\sqrt[3]{2^{7}\cdot5^{4}}$
  \task $\sqrt[4]{3^{9}}$
  \task $\sqrt{a^{7}\,b^{4}}$
  \task $\sqrt[3]{x^{5}\,y^{9}}$
\end{tasks}

\ex{3} Extrae todos los factores posibles:
\begin{tasks}(4)
  \task $\sqrt{\dfrac{50}{81}}$
  \task $\sqrt{12a^{3}}$
  \task $\sqrt[3]{-54}$
  \task $\sqrt[4]{32\,a^{6}\,b^{5}}$
  \task $\sqrt{\dfrac{72}{25}}$
  \task $\sqrt[3]{-16a^{4}}$
  \task $\sqrt{18\,a^{5}\,b^{2}}$
  \task $\sqrt[3]{81\,x^{7}\,y^{3}}$
\end{tasks}

\ex{1} Introduce los factores dentro del radical:
\begin{tasks}(6)
  \task $2\sqrt{5}$
  \task $3\sqrt{3}$
  \task $5\sqrt{2}$
  \task $2\sqrt[3]{3}$
  \task $3\sqrt[4]{2}$
  \task $a\sqrt{b}$
  \task $2^{3}\sqrt{2}$
  \task $3^{2}\sqrt{5}$
  \task $2^{2}\sqrt[3]{2}$
  \task $5^{2}\sqrt[3]{5^{2}}$
  \task $2^{3}\cdot3\sqrt{3}$
  \task $a^{2}\sqrt[4]{a^{3}}$
\end{tasks}

\ex{2} Introduce los factores dentro del radical y ordena de menor a mayor:
\[
3\sqrt{5}\qquad 4\sqrt{3}\qquad 2\sqrt{11}\qquad 5\sqrt{2}
\]

% =====================================================
\seccio{6}{Operaciones con radicales}

\ex{1} Opera estos radicales semejantes:
\begin{tasks}(2)
  \task $3\sqrt{5}+2\sqrt{5}$
  \task $7\sqrt{3}-4\sqrt{3}+\sqrt{3}$
  \task $2\sqrt[3]{6}-5\sqrt[3]{6}+8\sqrt[3]{6}$
  \task $\sqrt{2}+3\sqrt{2}-6\sqrt{2}$
\end{tasks}

\ex{2} Extrae factores y opera:
\begin{tasks}(2)
  \task $\sqrt{8}+\sqrt{18}$
  \task $\sqrt{12}+\sqrt{75}-\sqrt{27}$
  \task $2\sqrt{20}-\sqrt{45}+\sqrt{80}$
  \task $\sqrt{50}-2\sqrt{32}+3\sqrt{2}$
  \task $\sqrt[3]{16}+\sqrt[3]{54}-\sqrt[3]{2}$
  \task $\sqrt{28}+\sqrt{63}-\sqrt{7}$
\end{tasks}

\ex{1} Calcula y simplifica:
\begin{tasks}(6)
  \task $\sqrt{3}\cdot\sqrt{12}$
  \task $\sqrt{5}\cdot\sqrt{20}$
  \task $\sqrt{2}\cdot\sqrt{6}$
  \task $\sqrt[3]{2}\cdot\sqrt[3]{4}$
  \task $\sqrt{48}:\sqrt{3}$
  \task $\sqrt[3]{54}:\sqrt[3]{2}$
\end{tasks}

\ex{2} Calcula y simplifica:
\begin{tasks}(3)
  \task $2\sqrt{3}\cdot5\sqrt{6}$
  \task $\sqrt{10}\cdot\sqrt{15}$
  \task $\left(3\sqrt{2}\right)^{2}$
  \task $12\sqrt{30}:\left(4\sqrt{6}\right)$
  \task $\left(\sqrt{3}\right)^{5}$
  \task $\sqrt{2}\cdot\sqrt{3}\cdot\sqrt{6}$
\end{tasks}

\ex{3} Opera y simplifica:
\begin{tasks}(3)
  \task $\sqrt{2}\left(\sqrt{8}+\sqrt{18}\right)$
  \task $\sqrt{3}\left(2\sqrt{3}-\sqrt{12}\right)$
  \task $\left(\sqrt{6}+\sqrt{2}\right)\cdot\sqrt{2}$
  \task $\sqrt{12}+\sqrt{27}-\sqrt{3}\cdot\sqrt{6}$
  \task $\dfrac{\sqrt{50}-\sqrt{8}}{\sqrt{2}}$
\end{tasks}

% =====================================================
\seccio{7}{Intervalos: definición y tipos}

\ex{1} Completa la tabla:\par\smallskip
{\renewcommand{\arraystretch}{1.75}\centering
\begin{tabular}{|C{3.6cm}|C{4.2cm}|C{6.2cm}|}
\hline
\rowcolor{soft}\textbf{Intervalo} & \textbf{Desigualdad} & \textbf{Tipo de intervalo} \\ \hline
$\lo -3,\ 5\ro $ & & \\ \hline
$[-1,\ 4]$ & & \\ \hline
 & $2<x\le 6$ & \\ \hline
 & $x\ge -2$ & \\ \hline
$\lo -\infty,\ 0\ro $ & & \\ \hline
 & $-5\le x<-1$ & \\ \hline
\end{tabular}\par}

\ex{1} Escribe como intervalo y como desigualdad:
\begin{tasks}(2)
  \task Los números mayores que $-4$.
  \task Los números menores o iguales que $7$.
  \task Los números comprendidos entre $0$ y $10$, ambos incluidos.
  \task Los números mayores que $-2$ y menores o iguales que $3$.
\end{tasks}

\ex{2} Indica si el número pertenece o no al intervalo:
\begin{tasks}(3)
  \task ¿$3\in\lo -1,\ 3\ro $?
  \task ¿$-2\in[-2,\ 5\ro $?
  \task ¿$0\in\lo -\infty,\ 0]$?
  \task ¿$2,5\in\lo 2,\ 3\ro $?
  \task ¿$-7\in\lo -6,\ +\infty\ro $?
  \task ¿$\sqrt{2}\in[1,\ 2]$?
\end{tasks}

\ex{3} Responde:
\begin{tasks}(1)
  \task Escribe tres números que pertenezcan al intervalo $\lo 1,\ 2\ro $. ¿Cuántos números reales contiene este intervalo?
  \task Escribe un intervalo que contenga el $0$ y el $5$, pero no el $-1$ ni el $6$.
  \task ¿Cuál es el mayor número entero que pertenece a $[-3,\ 4\ro $? ¿Y el mayor que pertenece a $\lo -3,\ 4]$?
  \task Escribe como intervalo el conjunto de números reales cuya distancia al $2$ es menor que $3$.
\end{tasks}

% =====================================================
\seccio{8}{Representación gráfica de intervalos}

\ex{1} Representa en la recta real los siguientes intervalos:\par\smallskip
\begin{tabular}{@{}l@{\ \ }l@{\qquad}l@{\ \ }l@{}}
a) $[-2,\ 3]$ & \rectabuida & b) $\lo 0,\ 4\ro $ & \rectabuida \\[16pt]
c) $\lo -\infty,\ 1]$ & \rectabuida & d) $\lo -3,\ +\infty\ro $ & \rectabuida \\[16pt]
e) $[-4,\ -1\ro $ & \rectabuida & f) $x\le 2$ & \rectabuida
\end{tabular}

\ex{2} Escribe el intervalo y la desigualdad que corresponden a cada representación:\par\smallskip
\begin{tabular}{@{}l@{\ \ }l@{\qquad}l@{\ \ }l@{}}
a) & \rectaint{-1}{4}{b}{p} & b) & \rectaint{2}{6.4}{p}{i} \\[16pt]
c) & \rectaint{-6.4}{-3}{i}{b} & d) & \rectaint{-4}{0}{p}{p}
\end{tabular}

\ex{3} Representa en la recta real y escribe como intervalo (o unión de intervalos):\par\smallskip
\begin{tabular}{@{}l@{\ \ }l@{}}
a) $x>-1$\ \ y\ \ $x\le 5$ & \rectabuida \\[16pt]
b) $x\ge -3$\ \ y\ \ $x<2$ & \rectabuida \\[16pt]
c) $x\le -1$\ \ o\ \ $x>4$ & \rectabuida
\end{tabular}

% =====================================================
\seccio{9}{Unión e intersección de intervalos}

\ex{1} Sean $A=[-2,\ 4]$ y $B=\lo 1,\ 6\ro $. Represéntalos en la misma recta y calcula $A\cup B$ y $A\cap B$.\par\smallskip
{\centering\rectabuida\par\vspace{12pt}\rectabuida\par}

\ex{2} Calcula $A\cup B$ y $A\cap B$ en cada caso (haz antes su representación gráfica):
\begin{tasks}(2)
  \task $A=\lo -\infty,\ 3\ro $,\ \ $B=[0,\ +\infty\ro $
  \task $A=[-5,\ -1]$,\ \ $B=\lo -3,\ 2\ro $
  \task $A=\lo -\infty,\ -2]$,\ \ $B=\lo 1,\ +\infty\ro $
  \task $A=[0,\ 5\ro $,\ \ $B=[5,\ 8]$
  \task $A=\lo -4,\ 4\ro $,\ \ $B=[-1,\ 1]$
\end{tasks}

\ex{3} Responde:
\begin{tasks}(1)
  \task Escribe dos intervalos $A$ y $B$ tales que $A\cap B=[2,\ 3]$.
  \task Sean $A=[-3,\ 5\ro $ y $B=\lo a,\ 8]$. ¿Qué valor tiene que tener $a$ para que $A\cap B=\lo 1,\ 5\ro $?
  \task Sean $A=[-2,\ 6]$, $B=\lo 0,\ +\infty\ro $ y $C=\lo -\infty,\ 3]$. Calcula $A\cap B\cap C$ y $(A\cap B)\cup C$.
  \task Escribe el intervalo $[-1,\ 7\ro $ como la intersección de dos semirrectas.
\end{tasks}

% =====================================================
\seccio{10}{Clasificación de los números reales}

\ex{1} Marca con una cruz los conjuntos a los que pertenece cada número:\par\smallskip
{\renewcommand{\arraystretch}{1.3}\centering
\begin{tabular}{|C{3.6cm}|C{1.3cm}|C{1.3cm}|C{1.3cm}|C{1.3cm}|C{1.3cm}|}
\hline
\rowcolor{soft}\textbf{Número} & $\mathbb{N}$ & $\mathbb{Z}$ & $\mathbb{Q}$ & $\mathbb{I}$ & $\mathbb{R}$ \\ \hline
$12$ & & & & & \\ \hline
$-5$ & & & & & \\ \hline
$\dfrac37$ & & & & & \\ \hline
$0,25$ & & & & & \\ \hline
$\sqrt{3}$ & & & & & \\ \hline
$\pi$ & & & & & \\ \hline
$-\sqrt{49}$ & & & & & \\ \hline
$0,\overline{3}$ & & & & & \\ \hline
$1,234\,567\,891\,011\ldots$ & & & & & \\ \hline
\end{tabular}\par}

\ex{2} Indica el conjunto numérico más pequeño al que pertenece cada número:
\begin{tasks}(4)
  \task $\dfrac{18}{6}$
  \task $-\sqrt{25}$
  \task $\sqrt{8}$
  \task $2,\overline{5}$
  \task $1+\sqrt{2}$
  \task $-\dfrac{27}{9}$
  \task $\sqrt{0,09}$
  \task $0,101\,001\,000\,1\ldots$
  \task $\sqrt[3]{-64}$
  \task $\dfrac{\pi}{2}$
\end{tasks}

\ex{2} Indica si las siguientes afirmaciones son verdaderas o falsas. Justifícalo.
\begin{tasks}(2)
  \task Todo número natural es entero.
  \task Todo número entero es natural.
  \task Todo número decimal periódico es racional.
  \task La raíz cuadrada de un natural siempre es irracional.
  \task Hay números que son racionales e irracionales a la vez.
  \task Todos los números decimales son racionales.
\end{tasks}

\ex{3} Responde:
\begin{tasks}(1)
  \task Escribe un número irracional comprendido entre $3$ y $4$.
  \task Escribe un número racional comprendido entre $\sqrt{2}$ y $\sqrt{3}$.
  \task ¿Para qué valores naturales de $n$, entre $1$ y $30$, el número $\sqrt{n}$ es racional?
  \task ¿La suma de dos números irracionales es siempre irracional? Busca un contraejemplo.
\end{tasks}

\vfill
\begin{center}\color{muted}\small\faCheckCircle[regular]\ \ ¡Ánimo! Si te cuesta algún ejercicio, repasa los apuntes interactivos de la unidad.\end{center}

\end{document}
